\section*{Hoofdstuk \ref{ch:g9:landforms} --- Landvormen}

\begin{solution}{exo:g9:landforms:1}
(a) De wind. (b) Water dat het \omterm{def:g3:stones-pebbles-sand:rock}{gesteente} oplost (en in de \omterm{def:g9:landforms:karst}{grot} calciet
afzet). (c) De \omterm{def:g2:water-on-earth:ocean}{zee} (\omterm{def:g9:landforms:beach}{kustlangstransport}). (d) IJs. (e) De \omterm{def:g2:water-on-earth:ocean}{zee}. (f) Een
\omterm{def:g1:landscapes-around-us:water}{rivier}.
\end{solution}

\begin{solution}{exo:g9:landforms:2}
(a) \omterm{def:g9:landforms:uplift}{Opheffing}. (b) De \omterm{def:g9:landforms:dune}{lijzijde}. (c) Een \omterm{def:g9:landforms:karst}{zinkgat}. (d) \omterm{def:g9:landforms:beach}{Kustlangstransport}.
(e) \omterm{def:g9:landforms:karst}{Karst}.
\end{solution}

\begin{solution}{exo:g9:landforms:3}
(a) Niet waar: de hoorns wijzen met de wind mee; de ronde neus is naar de
wind gekeerd. (b) Waar. (c) Niet waar: de wind verplaatst het korrel voor
korrel, de flauwe helling op en langs de \omterm{def:g9:landforms:dune}{lijzijde} naar beneden. (d) Waar.
\end{solution}

\begin{solution}{exo:g9:landforms:4}
(b), (d), (e), (a), (c).
\end{solution}

\begin{solution}{exo:g9:landforms:5}
(a) Uit het westen: de hoorns wijzen met de wind mee. (b) Aan de oostkant,
tussen de hoorns. (c) $60 \div 6 = 10$ jaar. (d) Een groter \omterm{def:g9:landforms:dune}{duin} bevat meer
\omterm{def:g3:stones-pebbles-sand:sand}{zand}, en al dat \omterm{def:g3:stones-pebbles-sand:sand}{zand} moet over de kam worden gevoerd voordat het \omterm{def:g9:landforms:dune}{duin}
vooruitkomt.
\end{solution}

\begin{solution}{exo:g9:landforms:6}
(a) $7 \times 1 = \qty{7}{m}$. (b) $2 \div 12 = \frac{1}{6}$ van een jaar.
(c) Zeven jaar terugwijken in $\frac{1}{6}$ van een jaar: $7 \times 6 = 42$
keer sneller. (d) De golven slijten eerst een inkeping uit; het \omterm{def:g3:stones-pebbles-sand:rock}{gesteente}
erboven hangt er een tijdlang overheen, zonder te bewegen; dan stort de
overhang in één keer in, en de rand springt terug.
\end{solution}

\begin{solution}{exo:g9:landforms:7}
(a) Naar het oosten. (b) Drie bewegingen het \omterm{def:g9:landforms:beach}{strand} op, schuin naar het
noordoosten, elk gevolgd door een beweging recht terug naar beneden, naar
het zuiden: een zigzag die naar het oosten opschuift. (c) Een \omterm{def:g9:landforms:beach}{landtong}, die
naar het oosten dwars over de monding van de baai groeit. (d) De \omterm{def:g9:landforms:beach}{landtong}
beschut het water erachter tegen de golven; in rustig water kan de fijne
modder bezinken.
\end{solution}

\begin{solution}{exo:g9:landforms:8}
(a) Het neemt kooldioxide op uit de lucht en uit de bodem. (b) Het zwakke
zuur lost kalksteen langzaam op, maar tast graniet nauwelijks aan. (c) Het
water kan het \omterm{def:g3:stones-pebbles-sand:rock}{gesteente} alleen binnendringen via de spleten en langs de
lagen; het verwijdt eerst die openingen. (d) Het \omterm{def:g3:stones-pebbles-sand:rock}{gesteente} onder het
oppervlak lost op en de grond zakt langzaam in; of het dak van een \omterm{def:g9:landforms:karst}{grot}
stort in.
\end{solution}

\begin{solution}{exo:g9:landforms:9}
(a) Als de druppel de lucht van de \omterm{def:g9:landforms:karst}{grot} raakt, staat hij een deel van zijn
kooldioxide af en kan hij niet meer alle opgeloste kalksteen vasthouden,
die als calciet neerslaat. (b) Een stalagmiet groeit uit druppels die op de
bodem uiteenspatten en zich verspreiden; er hangt geen druppel aan zijn
punt die een buis kan bouwen.
(c) $\qty{30}{cm} = \qty{300}{mm}$; $300 \div 0.1 = 3000 = 3 \times 10^3$
jaar.
\end{solution}

\begin{solution}{exo:g9:landforms:10}
(a) $4600 \div 10^7 = 4.6 \times 10^{-4}$~m per jaar, of \qty{0.46}{mm} per
jaar. (b) $0.46 - 0.2 = \qty{0.26}{mm}$ per jaar. (c) $\qty{1}{m} =
\qty{1000}{mm}$; $1000 \div 0.46 \approx 2200$, ongeveer tweeduizend keer
sneller. (d) Ze wordt miljoenen jaren lang volgehouden: $\qty{0.26}{mm} \times
10^7 = 2.6 \times 10^6$~mm, ongeveer \qty{2.6}{km}.
\end{solution}

\begin{solution}{exo:g9:landforms:11}
(a) Vorm: sikkelduinen met hun \omterm{def:g9:landforms:dune}{lijzijde} aan de zuidoostkant. Kracht en
proces: de wind, die uit het noordwesten waait; \omterm{def:g5:water-shapes-the-land:deposition}{afzetting}. Materiaal: los
\omterm{def:g3:stones-pebbles-sand:sand}{zand}. (b) Vorm: een \omterm{def:g9:landforms:cliff}{klif} boven een vlakke rotsbank die bij eb droogvalt.
Kracht en proces: de \omterm{def:g2:water-on-earth:ocean}{zee}, \omterm{def:g5:water-shapes-the-land:erosion}{erosie}; de rotsbank is een abrasieplatform dat
het terugwijkende \omterm{def:g9:landforms:cliff}{klif} heeft achtergelaten. Materiaal: wit \omterm{def:g3:stones-pebbles-sand:rock}{gesteente},
waarschijnlijk krijt, een zacht \omterm{def:g3:stones-pebbles-sand:rock}{gesteente}. (c) Vorm: een beek die in een
gat verdwijnt, een \omterm{def:g6:water-underground:spring}{bron} ver daarvandaan. Kracht en proces: water dat het
\omterm{def:g3:stones-pebbles-sand:rock}{gesteente} oplost en onder de grond stroomt. Materiaal: kalksteen; het
plateau is een karstgebied.
\end{solution}

\begin{solution}{exo:g9:landforms:12}
(a) Het kleine \omterm{def:g9:landforms:dune}{duin}: $45 \div 6 = 7.5$ jaar; het grotere: $45 \div 3 =
15$ jaar. (b) $45 \div 100 = 0.45$ van een jaar, ongeveer vijf à zes
maanden. (c) Plant bijvoorbeeld grassen op de \omterm{def:g9:landforms:dune}{duinen}: hun wortels houden
het \omterm{def:g3:stones-pebbles-sand:sand}{zand} vast, zodat de wind het niet meer kan meenemen.
\end{solution}

\begin{solution}{pb:g9:landforms:1}
\textbf{1.} $46.0 - 30.4 = \qty{15.6}{m}$.

\textbf{2.} $15.6 \div 20 = \qty{0.78}{m}$ per jaar.

\textbf{3.} Tussen 2012 en 2014: $41.0 - 36.2 = \qty{4.8}{m}$ in twee
jaar, of $4.8 \div 2 = \qty{2.4}{m}$ per jaar.

\textbf{4.} $2.4 \div 0.78 \approx 3$ keer sneller. Waarschijnlijk een
stormachtige winter, waarin grote overhangen zijn ingestort.

\textbf{5.} \qty{78}{cm} per jaar; $0.78 \times 100 = \qty{78}{m}$ per
eeuw.

\textbf{6.} Ja: \qty{0.78}{m} per jaar is iets minder dan een meter per jaar.

\textbf{7.} $17 \div 0.78 \approx 22$ jaar.

\textbf{8.} Het \omterm{def:g9:landforms:cliff}{klif} wijkt met sprongen terug: een periode van twee jaar kan
een grote instorting bevatten, of helemaal geen. Over 20 jaar middelen de
sprongen en de rustige jaren zich uit.

\textbf{9.} De rand moet nog $30.4 - 10 = \qty{20.4}{m}$ afleggen: $20.4 \div 0.78
\approx 26$ jaar, dus in 2050.

\textbf{10.} $2 \times 0.78 = \qty{1.56}{m}$ per jaar; $30.4 \div 1.56
\approx 19.5$ jaar na 2024, dus in 2043.

\textbf{11.} De rand van een \omterm{def:g9:landforms:cliff}{klif} kan op elk moment afbrokkelen; erop
staan, of er dichtbij, is gevaarlijk.

\textbf{12.} $30.4 \div 0.78 \approx 39$ jaar na 2024: in
\textbf{2063}.
\end{solution}
